\documentclass[11pt]{article}
\usepackage[margin=1in]{geometry}
\usepackage{amsmath,amssymb,amsthm,mathtools}
\usepackage{graphicx}
\usepackage[hidelinks]{hyperref}
\usepackage{xcolor}

\newtheorem{theorem}{Theorem}
\newtheorem{proposition}[theorem]{Proposition}
\newtheorem{remark}[theorem]{Remark}
\newcommand{\Sp}{\mathcal S_p}
\newcommand{\dd}{\,\mathrm d}

\title{The mixed-kernel Schatten estimate fails below $p=2$\\
\large A counterexample inside the fractional-commutator class}
\author{Autonomous mathematical discovery run \texttt{fa\_banach\_001}}
\date{17 August 2026}

\begin{document}
\maketitle

\noindent\textbf{Status.}
Candidate counterexample, likely valid pending expert review.

\section{Source question}

Gerontogiannis and Mesland consider an Ahlfors $d_f$-regular compact metric
measure space $(X,d,\mu)$ and, for $0<2\alpha<\beta\leq1$, the commutator
kernel
\begin{equation}\label{eq:kernel}
 K_{\alpha,h}(x,y)=\frac{h(y)-h(x)}{d(x,y)^{d_f+2\alpha}},
 \qquad h\in\operatorname{H\ddot ol}_{\beta}(X,d).
\end{equation}
In Remark 2.18 of \cite{GM}, on PDF page 20, they ask whether for
$1<p<2$ and $q=p/(p-1)$ the estimate
\begin{equation}\label{eq:source}
 \|K_{\alpha,h}\|_{\mathcal S_p}
 \leq
 \left(\int_X\left(\int_X|K_{\alpha,h}(x,y)|^q\dd\mu(x)
 \right)^{p/q}\dd\mu(y)\right)^{1/p}
\end{equation}
continues to hold.  For $p\geq2$ this is the symmetric-kernel specialization
of the Russo--Goffeng mixed-norm estimate \cite{Goffeng}.  An affirmative
answer would improve the threshold in their Proposition 2.17.

\begin{center}
\includegraphics[width=0.96\textwidth]{figures/open_problem_crop.png}
\end{center}

\section{Proof intuition}

The binary ultrametric Cantor space allows the infinite-dimensional kernel to
be reduced exactly to four cylinders.  Choose $h$ to have two very close
values on each of the two first-level halves.  The larger fractional weight
inside a half compensates for that small difference.  As the ultrametric
scale tends to zero, all six entries above the diagonal of the resulting
$4\times4$ antisymmetric matrix tend to one.

That limiting matrix has singular values $\sqrt2+1$ and $\sqrt2-1$, each with
multiplicity two, whereas every column has three unit entries.  The resulting
scalar inequality is strict in the wrong direction for every $1<p<2$.
Continuity then gives a genuine example at a finite ultrametric scale.  This
directly addresses the restrictive reading in which the source's $K$ must be
the fractional commutator kernel \eqref{eq:kernel}.

\section{Counterexample}

\begin{theorem}\label{thm:main}
For every $1<p<2$, there exist
\begin{itemize}
\item a compact Ahlfors-regular metric measure space $(X,d,\mu)$,
\item parameters $\alpha=1/8$ and $\beta=1$, and
\item a real locally constant (hence Lipschitz) function $h$,
\end{itemize}
such that $p>d_f/(\beta-2\alpha)$ and the fractional commutator kernel
\eqref{eq:kernel} violates \eqref{eq:source}.  In particular, the answer to
Remark 2.18, with the displayed constant-one inequality as written, is
negative throughout $1<p<2$.
\end{theorem}

\begin{proof}
Fix $p\in(1,2)$ and write $q=p/(p-1)$.  Let
$X=\{0,1\}^{\mathbb N}$ with Bernoulli measure assigning mass $2^{-k}$ to
each cylinder of length $k$.  For $0<r<1$, set
\[
 d_r(x,y)=r^{j(x,y)},
\]
where $j(x,y)\geq0$ is the index of the first coordinate at which $x$ and
$y$ differ.  This is the standard binary ultrametric, and $\mu$ is Ahlfors
regular of dimension
\[
 d_f=\frac{\log2}{\log(1/r)}.
\]
Indeed, balls at the discrete radii $r^k$ are cylinders of measure $2^{-k}$;
the estimates between adjacent radii change only by fixed factors.

Set $\alpha=1/8$ and $\beta=1$.  Denote the four length-two cylinders, in
lexicographic order, by $E_0,E_1,E_2,E_3$, and put
\[
 c=r^{-(d_f+2\alpha)}=2r^{-1/4}.
\]
Define the locally constant real function $h$ by the values
\begin{equation}\label{eq:hvalues}
 (h|_{E_0},h|_{E_1},h|_{E_2},h|_{E_3})
   =(0,c^{-1},1,1+c^{-1}).
\end{equation}
This function is Lipschitz: it is constant within every length-two cylinder,
and only finitely many positive distances occur between distinct such
cylinders.

The kernel $K_{\alpha,h}$ vanishes on every $E_i\times E_i$.  Between two
distinct length-two cylinders it is constant.  The distance between sibling
cylinders is $r$, while the distance between cylinders in different
first-level halves is one.  Consequently, on the orthonormal basis
$e_j=2\mathbf1_{E_j}$ of the four-dimensional cylinder-constant subspace,
the integral operator has matrix $A_c/4$, where $A_c$ is real antisymmetric
and its upper-triangular entries are
\begin{equation}\label{eq:Ac}
 (A_{01},A_{02},A_{03},A_{12},A_{13},A_{23})
  =(1,1,1+c^{-1},1-c^{-1},1,1).
\end{equation}
The operator annihilates the orthogonal complement of this subspace, because
the kernel is constant in each variable on every cylinder pair.  Thus its
nonzero singular values are exactly those of $A_c/4$.

As $r\downarrow0$, we have $c\to\infty$ and $A_c\to A$, where
\begin{equation}\label{eq:A}
 A=\begin{pmatrix}
 0&1&1&1\\
 -1&0&1&1\\
 -1&-1&0&1\\
 -1&-1&-1&0
 \end{pmatrix}.
\end{equation}
The singular values of $A$ are
\[
 s_+=\sqrt2+1,\quad s_+=\sqrt2+1,\quad
 s_-=\sqrt2-1,\quad s_-=\sqrt2-1.
\]
For completeness, $A$ is skew-symmetric, its squared Frobenius norm is $12$,
and its Pfaffian is $1$.  Hence the two distinct squared singular values have
sum $6$ and product $1$, giving $3\pm2\sqrt2$.

Every column of $A$ has three entries of modulus one.  Therefore, for the
limiting four-block kernel, the $p$th power of the ratio of the left side of
\eqref{eq:source} to the right side is
\begin{equation}\label{eq:ratio}
 R_p^p=
 \frac{(\sqrt2+1)^p+(\sqrt2-1)^p}{2\,3^{p-1}}.
\end{equation}
We show that $R_p>1$ whenever $1<p<2$.  Put $a=1+\sqrt2$, so that
$a^{-1}=\sqrt2-1$, and define
\[
 F(t)=\log\!\left(\frac{a^t+a^{-t}}2\right)-(t-1)\log3.
\]
Since $a^2+a^{-2}=6$, one has $F(2)=0$.  Moreover,
\[
 F'(t)=\log(a)\frac{a^t-a^{-t}}{a^t+a^{-t}}-\log3
 <\log(a)-\log3<0.
\]
Thus $F$ is strictly decreasing and $F(p)>F(2)=0$ for $p<2$.  Equation
\eqref{eq:ratio} gives $R_p>1$.

Both sides of \eqref{eq:source} are continuous functions of the entries of
the finite matrix $A_c$.  Hence, for all sufficiently large finite $c$
(equivalently, all sufficiently small positive $r$), $A_c$ still violates
\eqref{eq:source}.  Finally,
\[
 \frac{d_f}{\beta-2\alpha}
 =\frac{4\log2}{3\log(1/r)}\longrightarrow0
 \qquad(r\downarrow0),
\]
so the same choice of sufficiently small $r$ also ensures
$p>d_f/(\beta-2\alpha)$.  All hypotheses relevant to the proposed improved
range are satisfied.
\end{proof}

\section{A sharp obstruction for arbitrary antisymmetric kernels}

The preceding theorem refutes the constant-one estimate inside the special
fractional family.  The next observation explains why no generic
antisymmetric-kernel argument can repair it by inserting a constant depending
only on $p$.

\begin{proposition}\label{prop:unbounded}
Fix $1<p<2$.  Among bounded real antisymmetric step kernels on the Ahlfors
regular space $[0,1]$, the ratio of the two sides of \eqref{eq:source} is
unbounded.  More precisely, for every power of two $n$ there is such a kernel
with ratio exactly $n^{1/p-1/2}$.
\end{proposition}

\begin{proof}
Let $H_n$ be the Sylvester Hadamard matrix and $U_n=n^{-1/2}H_n$, so $U_n$ is
real orthogonal and all its entries have modulus $n^{-1/2}$.  Set
\[
 B_n=\begin{pmatrix}0&U_n\\-U_n^{\mathsf T}&0\end{pmatrix}.
\]
Then $B_n$ is a real antisymmetric orthogonal matrix of size $2n$.  Partition
$[0,1]$ into $2n$ equal intervals and use $B_n$ as the matrix of a step
kernel.  The integral operator has matrix $B_n/(2n)$ on normalized interval
indicators, hence Schatten norm $(2n)^{1/p-1}$.

Each kernel column is nonzero on exactly $n$ intervals, with modulus
$n^{-1/2}$.  Its mixed norm is
\[
 2^{-1/q}n^{-1/2}.
\]
The quotient is therefore
\[
 \frac{(2n)^{1/p-1}}{2^{-1/q}n^{-1/2}}
 =n^{1/p-1/2}\longrightarrow\infty.
\]
\end{proof}

\section{Scope, verification, and novelty}

Theorem \ref{thm:main} disproves the exact inequality (2.12) with constant
one, even when $K$ is the source's own fractional commutator kernel.  It does
\emph{not} rule out an estimate with a factor depending on $p$ (and possibly
on geometric regularity constants), nor does it rule out proving the
$\ell=1$ improvement by a different argument.  Proposition
\ref{prop:unbounded} shows only that such a repair cannot hold uniformly for
all antisymmetric kernels; the rigid fractional-difference structure would
have to be used.

The accompanying script checks the singular values, finite-$c$ ratios, and
the exact Hadamard family formulas.  These checks are not part of the proof.

A bounded search on 17 August 2026 covered the current arXiv record for
\cite{GM} (still v1), exact queries for Remark 2.18 and inequality (2.12), and
fractional-commutator, ultrametric, and mixed Schatten-norm variants.  No later
resolution or this four-cylinder construction was found.  The classical fact
that kernel criteria below $p=2$ are substantially more delicate, including
Carleman-type obstructions, is reviewed by Delgado and Ruzhansky
\cite{DR}; that general literature does not supply the special
fractional-commutator example above.  Novelty confidence is moderate pending
specialist review.

\paragraph{Human review recommendation.}
Check the cylinder-distance convention and the resulting six entries in
\eqref{eq:Ac}; the finite-rank reduction; and whether the authors intended an
unspecified multiplicative constant rather than the displayed constant-one
inequality.  The last point affects the scope, not the validity, of the stated
counterexample.

\begin{thebibliography}{9}
\bibitem{GM}
D. M. Gerontogiannis and B. Mesland,
\emph{Ideal quantum metrics from fractional Laplacians},
arXiv:2502.04187v1 (2025).

\bibitem{Goffeng}
M. Goffeng,
\emph{Analytic formulas for the topological degree of non-smooth mappings:
the odd-dimensional case},
Adv. Math. 231 (2012), 357--377.

\bibitem{DR}
J. Delgado and M. Ruzhansky,
\emph{Schatten--von Neumann classes of integral operators},
J. Math. Pures Appl. 154 (2021), 1--29.
\end{thebibliography}

\end{document}
